\documentclass[a4paper]{article}
\usepackage{graphicx} % Required for inserting images
\usepackage{wrapfig}
\usepackage{enumitem}
\usepackage[center]{caption}
\usepackage[english, russian]{babel}
\usepackage{amsmath}
\usepackage{geometry}
\geometry{verbose, a4paper, tmargin=2cm, bmargin=2cm, lmargin=2.0cm, rmargin=2.0cm}

\begin{document}

\begin{titlepage}
	\centering
	{\scshape\Large МОСКОВСКИЙ ФИЗИКО-ТЕХНИЧЕСКИЙ ИНСТИТУТ \\
	(НАЦИОНАЛЬНЫЙ ИССЛЕДОВАТЕЛЬСКИЙ УНИВЕРСИТЕТ)\\ % название ВУЗа большим шрифтом
    Физтех-школа аэрокосмических технологий}
	
	\vspace{4cm} % отступ 4 см по вертикали
	{\LARGE Отчет о выполнении лабораторной работы}
	
	\vspace{1cm} % ещё 1 см
	{\huge\bf Изучение колебаний связанных маятников}
	
	\vspace{1cm} % ещё 1 с
	\vfill % заполнение по вертикали пустотой (LaTeX сам решит, сколько здесь отступить)
	
    \begin{flushright} % Этот текст будет с правого края страницы
	   {\LARGE Ефремова Татьяна, Б03-503}
    \end{flushright}
    \vfill

	\today % Дата сборки документа
\end{titlepage}

\section{Аннотация}
Цели работы: изучение гармонических (нормальных) колебаний связанной системы двух маятников. Изучение биений как результата сложения нормальных колебаний свзяанной системы двух маятников.

\section{Теоретические сведения}
Связанная колебательная система -- это совокупность нескольких маятников (парциальных систем), каким-либо образом связанных между собой. В данной работе ее представляют два математических маятника, скрепленных пружинами. \\
\begin{wrapfigure}{R}{0.3\textwidth}
    \includegraphics[width=1\linewidth]{C:/Users/User/Documents/mipt_takonef/лабы/ВПВ/scheme.JPG}
    \caption{Схема системы}
\end{wrapfigure}

\noindent Уравнение движения для каждого ма ятника можно получить из общего уравнения динамики вращательного движения:
\begin{equation}
J\ddot{\varphi} = M,
\end{equation}
где $J$ -- момент инерции тела относительно оси вращения, $M$ -- момент действующих на тело сил.\\\\
Для малых колебаний ($sin{\varphi} \approx \varphi$) применительно к каждому маятнику уравнение (1) имеет вид:
\[
m\ddot{\varphi_1} + mgl\varphi_1 = k(d_1^2 + d_2^2)(\varphi_2-\varphi_1), 
\]
\begin{equation}
m\ddot{\varphi_2} + mgl\varphi_2 = k(d_1^2 + d_2^2)(\varphi_1-\varphi_2).
\end{equation}
Предполагается, что момент силы трения много меньше $M$, т. е. затуханием можно принебречь.\\ \\
Общее решение системы (2) имеет вид:
\[
\varphi_1 = Acos(\omega_1t + \alpha_1) + Bcos(\omega_2t + \alpha_2),
\]
\begin{equation}
\varphi_2 = Acos(\omega_1t + \alpha_1) - Bcos(\omega_2t + \alpha_2),
\end{equation}
где $\omega_1 = \sqrt{\frac{g}{l}}$, $\omega_2 = \sqrt{\frac{g}{l} + 2\frac{k(d_1^2 + d_2^2)}{ml^2}}$, а постоянные $A, B, \alpha_1$ и $\alpha_2$ определяются из начальных условий. \\ \\
В общем случае каждый маятник осуществляет сложное колебание, которое не является гармоническим, поскольку складывается из двух независимых колебаний с частотами $\omega_1$ и $\omega_2$. Однако, специальным подбором начальных условий можно добиться возбуждения колебаний состемы с частотой, одинаковой для обоих маятников. Такие колебания называются \textit{нормальными}, а число их равно числу степеней свободы системы. В рассматриваемом случае нормальные колебания могут быть вызваны двумя способами:  \\ \\
\textit{Синфазные колебания} -- колебания маятников, происходящие в одной фазе с меньшей из нормальных частот $\omega_1$. Они описываются выражениями (3) при $B=0$, что соответсвует начальным условиям:
\begin{equation}
t=0, \quad \varphi_1=\varphi_2=\varphi_0, \quad \dot{\varphi_1}=\dot{\varphi_2},
\end{equation}
т. е. в начальный момент времени маятники имеют одинаковые скорости и углы отклонения. Роль связи в таком случае исчезает. \\ \\
\textit{Антифазные колебания} -- колебания маятников, происходящие в противофазе с большей из нормальных частот $\omega_2$. Они описываются выражениями (3) при $A=0$, что соответсвует начальным условиям:
\begin{equation}
t=0, \quad \varphi_1=-\varphi_2, \quad \dot{\varphi_1}=-\dot{\varphi_2},
\end{equation}
т. е. в начальный момент времени маятники имеют противоположные скорости и углы отклонения. \\ \\
В случае произвольных начальных условий будут одновременно совершаться оба нормальных колебания. При этом возникают \textit{биения} -- периодческое возрастание и убывание амплитуды колебаний каждого из маятников. Их удобно отслеживать при начальных условиях:
\begin{equation}
t=0, \quad \varphi_1=\varphi_0, \quad \varphi_2 = 0, \quad \dot{\varphi_1}=\dot{\varphi_2}=0,
\end{equation} 
При этом постоянные системы (3) имеют значения
\[
A=B=\frac{\varphi_0}{2}, \quad \alpha_1=\alpha_2=0.
\] \\\\
При их подстановке сиcтема (3) приобретает вид:
\[
\varphi_1 = \varphi_0cos(\frac{\omega_2-\omega_1}{2}t)cos(\frac{\omega_1+\omega_2}{2}t),
\]
\begin{equation}
\varphi_1 = \varphi_0sin(\frac{\omega_1-\omega_2}{2}t)sin(\frac{\omega_1+\omega_2}{2}t)
\end{equation}
Если связь между маятниками мала, то $\omega_2-\omega_1 \ll \omega_1+\omega_2$. \\ \\
В таком случае движение маятников можно считать колбеаниями, происходящими с частотой $\omega = \frac{\omega_1+\omega_2}{2}$ и амплитудами, медленно меняющимися по времени. Период изменений амплитуд выражается как
\begin{equation}
\tau = \frac{2\pi}{\omega_2-\omega_1}
\end{equation}
и называется \textit{периодом биений}. При этом фазы колебаний маятников отличаются на $\frac{\pi}{2}$. Уменьшение амплитуды одного из маятников приводит к увеличению амплитуды второго. Общая энергия системы (без учета потерь) остается постоянной. 

\begin{wrapfigure}{R}{0.4\textwidth}
    \includegraphics[width=1\linewidth]{C:/Users/User/Documents/mipt_takonef/лабы/ВПВ/stand.JPG}
    \caption{Установка}
\end{wrapfigure}

\section{Оборудование}

\subsection{Используемое оборудование}
Лабораторная установка показана на рис. 2. Маятники 1, 2 соединены пружинами 3, 4. Пружины закрепляются на фиксаторах, которые свободно перемещаются по резьбе стержней маятников. Частота свободных колебаний маятников регулируется перемещением грузов 7, 8. Начальные условие колебаний маятников задаются с помощью секторов, размещенных на штанге стартового устройства.  Они приводятся в движение поворотом рукоятки 9.  Сектора  10, 11 служат для запуска антифазных колебаний, 11, 12 -- для синфазных колебаний. На маятниках укреплены датчики угловой скорости 5, 6.  Передача данных на компьютер осуществляется при помощи USB-кабелей. \\ \\
К приборам и принадлежностям относится компьютер с программным обеспечением для регистрации и обработки данных.

\subsection{Инструментальные погрешности}
\vspace{0.15 cm}

\setdescription{leftmargin=0pt}
\begin{description}
    \item[линейка]:
\\ $\Delta_l=\pm 1 $ мм (цена деления).
    \item[система измерения времени]: 
\\ $\Delta_t=\pm0,005$ c (шаг квантования).
    \item[система измерения угловой скорости]: 
\\ $\Delta_\omega=\pm0,01$ рад/c (шаг квантования).
\end{description}
\newpage
\section{Результаты измерений и обработка данных}
Масса циллиндра поступательного маятника $M = 2900 \pm 5$ г. Длина нити подвеса $L = 2170 \pm 5$ мм.

\begin{table}[!h]
\centering
\caption{Амплитуды колебаний на шкале, скорости пуль}
\begin{tabular}{|c|c|c|c|}        \hline
    №Пули & m, мг & $\Delta x$, мм & $u$, м/с \\  \hline
    1 & $509$ & $11.6$ & $140$\\ \hline
    2 & $508$ & $11.5$ & $139$\\ \hline
    3 & $505$ & $11.4$ & $138$\\ \hline
    4 & $500$ & $11.4$ & $138$\\ \hline
    5 & $510$ & $11.5$ & $139$\\ \hline
\end{tabular}
\end{table}
\newpage

\noindent Рассчитаем начальную скорость пули по формуле:
\[u = \frac{M}{m} \sqrt{\frac{g}{L}} \Delta x\]
Рассчитаем погрешность скорости:

\[\varepsilon_{u_{\text{снст}}} = \varepsilon_M + \varepsilon_m + \frac{1}{2} \varepsilon_g + \frac{1}{2} \varepsilon_L + \varepsilon_{\Delta x} \approx 4.7\%\]

\[\sigma_{u_{\text{снст}}} = u \varepsilon_{u_{\text{снст}}} \approx 6.5 \text{м/с}\]

\[\sigma_{u_{\text{случ}}} \approx 5.6 \text{м/с}\]

\[\sigma_u = \sqrt{\sigma_{u_{\text{снст}}}^2 + \sigma_{u_{\text{случ}}}^2} \approx 8.6 \text{м/с} \]

\[\varepsilon_u = \frac{\sigma_u}{u} \approx 6.1\%\]
Тогда:

\[u = (139 \pm 9) \text{м/с}\]
\\
Измерим параметры крутильного баллистического маятника (рис. 2):
\begin{table}[!h]
\centering
\caption{Характеристики крутильного маятника}
\begin{tabular}{|c|c|c|c|c|}        \hline
    $d$, мм & $R$, мм & $r$, мм & $M_{g1}$, г& $M_{g2}$, г\\ \hline
    1387 & 336 & 213 & 725.6 & 724.5 \\ \hline
\end{tabular}
\end{table} \\
Результаты измерения времени одного полного колебания маятника (по результатам 10 колебаний):
\begin{table}[!h]
\centering
\caption{ Периоды колебаний крутильного маятника}
\begin{tabular}{|c|c|}        \hline
    Измерение & Период, с \\  \hline
    Без грузов, $T_1$ & $5.2\pm 0.1$  \\ \hline
    С грузами, $T_2$ & $6.9\pm 0.1$ \\ \hline
\end{tabular}
\end{table} \\
Отсюда:
\[\sqrt{kI} = \frac{4\pi MR^2 T_1}{T_1^2 - T_2^2} \approx (0.699 \pm 0.006)^{\frac{\text{кг} \cdot \text{м}^2}{\text{с}}}\]
Результаты измерения горизонтального смещения крутильного маятника на шкале, скорость пули для данной амплитуды.
\begin{table}[!h]
\centering
\caption{Амплитуды колебаний на шкале, скорости пуль}
\begin{tabular}{|c|c|c|c|}        \hline
    №Пули &  m, мг & $\Delta x$, мм & $u$, м/c \\  \hline
    6 & $510$ & $42$ & 99\\ \hline
    7 & $512$ & $39$ & 96\\ \hline
    8 & $509$ & $49$ & 102\\ \hline
    9 & $507$ & $45$ & 100\\ \hline
    10 & $513$ & $46$ & 102\\ \hline
\end{tabular}
\end{table} \\
Рассчитаем начальную скорость пули по формуле:

\[u = \varphi \frac{\sqrt{kI}}{mr} \approx x \frac{\sqrt{kI}}{2dmr}\]
Рассчитаем погрешность cкорости:
\[
\frac{\sigma_{80}}{g_0} = \sqrt{ \left( \frac{\sigma_L}{L} \right)^2 + 4 \left( \frac{\sigma_T}{T} \right)^2 }.
\]
\[\sigma_{u_{\text{случ}}} \approx 4^{\frac{\text{м}}{\text{с}}}\]

\[\sigma_{u_{\text{сист}}} = u \varepsilon_u \approx 6^{\frac{\text{м}}{\text{с}}}\]

\[\sigma_u = \sqrt{\sigma_{u_{\text{сист}}}^2 + \sigma_{u_{\text{случ}}}^2} \approx 7.2^{\frac{\text{м}}{\text{с}}}\]
По итогу:
\[u = (100 \pm 7)^{\frac{\text{м}}{\text{с}}}\] \\
Среднее значение и разброс: \\
Для поступательного маятника средняя скорость пули $u=139\pm 12$ м/с \\
Для крутильного маятника средняя скорость пули $u=100\pm 7$ м/с

\section{Выводы}
В резултьтате работы были измерены скорости пули для двух различных установок.

\end{document}
